%% ma: L11-B26-0007
%% lop: 11
%% chuong: 7
%% bai: 26
%% dang: 11.26.5
%% loai: TN4
%% muc: NB
%% dap_an: "A"
%% nguon: "(NBV)-ĐỀ SỐ 16-MỖI NGÀY 1 ĐỀ THI 2025.pdf (D16, MỖI NGÀY 1 ĐỀ - Đề số 16), Phần I Câu 4"
%% kien_thuc: "Lớp 11 Bài 26 (khoảng cách giữa hai đường thẳng song song)"
%% trang_thai: chinh-thuc
%% ghi_chu: "Đáp án gốc A đúng (lời giải gốc để trống). Hình lập phương dựng lại, nét khuất tính theo phép chiếu song song. Đã duyệt 17/10/2026."
\begin{ex}
Cho hình lập phương $ABCD.A'B'C'D'$ có cạnh bằng $a$. Khoảng cách giữa hai đường thẳng $AB$ và $DC$ bằng
\begin{center}
\begin{tikzpicture}[scale=.9]
  \coordinate (A) at (0,0);\coordinate (B) at (2.4,0);\coordinate (D) at (1,.7);\coordinate (C) at (3.4,.7);
  \coordinate (A') at (0,2.4);\coordinate (B') at (2.4,2.4);\coordinate (D') at (1,3.1);\coordinate (C') at (3.4,3.1);
  \draw (A)--(B)--(C)--(C')--(D')--(A')--cycle (A')--(B')--(C') (B)--(B');
  \draw[khuat] (A)--(D)--(C) (D)--(D');
  \foreach \p/\d in {A/below left,B/below right,C/right,D/above left,A'/left,B'/below right,C'/right,D'/above left}
    \fill (\p) circle (1pt) node[\d,font=\footnotesize] {$\p$};
\end{tikzpicture}
\end{center}
\choice{\True $a$}{$a\sqrt2$}{$a\sqrt3$}{$2a$}
\loigiai{$AB\parallel DC$ và $AD\perp AB$, $AD\perp DC$ nên $AD$ là đoạn vuông góc chung, do đó $d(AB,DC)=AD=a$. Chọn A.}
\end{ex}
